% !TeX program = xelatex
% ============================================================================
% 论文中文译文模板（双栏 / 中文期刊风格，参照高校学报双栏版式）
% 适用：华南理工大学学报、中文核心期刊一类的双栏投稿版式
% 用法：先用 scripts/probe_layout.py 量出版心与栏宽，改下面的参数
% ============================================================================
\documentclass[10pt,twocolumn]{ctexart}

% ---------- 版式参数 ----------
\newcommand{\PaperW}{210mm}
\newcommand{\PaperH}{297mm}
\newcommand{\MargL}{25mm}
\newcommand{\MargR}{25mm}
\newcommand{\MargT}{28mm}
\newcommand{\MargB}{25mm}
\newcommand{\ColSep}{7mm}      % 栏间距（实测两栏间空白）

\usepackage[
  paperwidth=\PaperW, paperheight=\PaperH,
  left=\MargL, right=\MargR, top=\MargT, bottom=\MargB,
  columnsep=\ColSep
]{geometry}

% ---------- 字体 ----------
\usepackage{fontspec}
\setmainfont{Times New Roman}
\setsansfont{Arial}
\setCJKmainfont{SimSun}
\setCJKsansfont{SimHei}
\setCJKfamilyfont{kai}{KaiTi}
\newcommand{\kai}{\CJKfamily{kai}}
\usepackage{unicode-math}
\setmathfont{STIX Two Math}

% ---------- 宏包 ----------
\usepackage{graphicx}
\usepackage{amsmath}
\usepackage{booktabs}
\usepackage{multirow}
\usepackage{array}
\usepackage{caption}
\usepackage{float}
\usepackage{enumitem}
\usepackage[hidelinks]{hyperref}

\captionsetup{font=small, labelsep=quad}
\setlist{nosep, leftmargin=1.5em}

% 双栏正文行距略紧，中文 10pt 用 1.22
\linespread{1.22}
\setlength{\parindent}{2em}
\setlength{\parskip}{0pt}

% ---------- 标题层级 ----------
\ctexset{
  section/format       = \raggedright\heiti\fontsize{11pt}{13pt}\selectfont,
  section/aftername    = \quad,
  section/beforeskip   = 10pt plus 2pt,
  section/afterskip    = 5pt,
  subsection/format    = \raggedright\heiti\fontsize{10pt}{12pt}\selectfont,
  subsection/aftername = \quad,
  subsection/beforeskip= 8pt,
  subsection/afterskip = 4pt,
  subsubsection/format = \raggedright\heiti\fontsize{10pt}{12pt}\selectfont,
  subsubsection/aftername = \quad,
}

% 页码
\pagestyle{plain}

\begin{document}

% ================= 通栏标题区 =================
\twocolumn[{%
\begin{center}
  {\heiti\fontsize{18pt}{22pt}\selectfont 中文标题（不超过 20 字为宜）\par}
  \vspace{6pt}
  {\fontsize{10.5pt}{13pt}\selectfont 作者一\textsuperscript{1}，作者二\textsuperscript{1}，作者三\textsuperscript{2}\par}
  \vspace{4pt}
  {\footnotesize\kai
   （1. 单位一，城市 邮编；2. 单位二，城市 邮编）\par}
  \vspace{8pt}
  \begin{minipage}{0.92\textwidth}
    \footnotesize
    \noindent\textbf{摘\quad 要：}摘要正文，楷体，约 200–300 字。\par
    \vspace{3pt}
    \noindent\textbf{关键词：}关键词一；关键词二；关键词三\par
    \vspace{3pt}
    \noindent\textbf{中图分类号：}TP301\quad\textbf{文献标志码：}A
  \end{minipage}
\end{center}
\vspace{6pt}
}]

% ================= 双栏正文 =================
\section{引言}

双栏正文。通栏图用 \texttt{figure*}，通栏表用 \texttt{table*}。

\begin{figure}[htbp]
  \centering
  \includegraphics[width=\linewidth]{figs/fig1.png}
  \caption{单栏图示例}
  \label{fig:one}
\end{figure}

\begin{table}[htbp]
  \centering
  \caption{单栏表示例}
  \label{tab:one}
  \footnotesize
  \begin{tabular}{@{}lrr@{}}
    \toprule
    参数 & 取值 & 单位 \\
    \midrule
    $Q$ & 200 & kg \\
    \bottomrule
  \end{tabular}
\end{table}

\input{refs}

\end{document}
